\section{Result 6: the smallest model}\label{sec:four}

You asked for width 4 to be kept in the record. I ran it on all four languages of the first
round and of the unequal-use round, three repeats each: 24 runs. By then many runs were sharing
the machine, so each went more than ten times slower than at the start of the night: in the 28
minutes (first-round languages) or 24 minutes (unequal use) that each had, it got through only
101 to 129 sets. One more run, on 16 animals, was trained for as long as the night allowed:
647 sets.

\textbf{Width 4 learned no facts.} Its score was 0 in 23 of the 24 runs and 0.3 in the other,
and it was still 0 after 647 sets in the long run. It did learn the shapes of the sentences. I
can tell from the finer measure of section \ref{sec:score}, how well a model predicts the next
word, on which lower is better. A model that knows how every sentence is built but no facts at
all scores 0.33 on it, and width 4 reached 0.36 by the end of the long run.

\textbf{How much these runs can bear.} In the first round, 112 to 129 sets are too few to
judge the two largest languages: width 8 had learned only part of them by then. Everywhere
else the comparison is fair. Width 8 knew every fact of the 16-animal language within 64 to
91 sets, and with unequal use it had a score of 77 or more on every language by set 91. Width 4 knew
nothing, and on 16 animals still nothing after 647 sets.

\textbf{This does not look like your idea at work.} Your idea says a model fails when it has
more to hold than fits. Width 4 fails with the fewest words. With 16 animals there are 128
bits to keep, and it has about 1,050 numbers in use (880 in its body and 43 rows of 4): about
6\% full by the rough guide. It seems to be too narrow to work at all. But I tried it for at
most 647 sets, and this night showed that a model can stand still for a long time and then
move, so even that is not proven.
